%-----------------------------------------------------------------------------------------------------------------------------------------------%
%	The MIT License (MIT)
%
%	Copyright (c) 2015 Jan Küster
%
%	Permission is hereby granted, free of charge, to any person obtaining a copy
%	of this software and associated documentation files (the "Software"), to deal
%	in the Software without restriction, including without limitation the rights
%	to use, copy, modify, merge, publish, distribute, sublicense, and/or sell
%	copies of the Software, and to permit persons to whom the Software is
%	furnished to do so, subject to the following conditions:
%
%	THE SOFTWARE IS PROVIDED "AS IS", WITHOUT WARRANTY OF ANY KIND, EXPRESS OR
%	IMPLIED, INCLUDING BUT NOT LIMITED TO THE WARRANTIES OF MERCHANTABILITY,
%	FITNESS FOR A PARTICULAR PURPOSE AND NONINFRINGEMENT. IN NO EVENT SHALL THE
%	AUTHORS OR COPYRIGHT HOLDERS BE LIABLE FOR ANY CLAIM, DAMAGES OR OTHER
%	LIABILITY, WHETHER IN AN ACTION OF CONTRACT, TORT OR OTHERWISE, ARISING FROM,
%	OUT OF OR IN CONNECTION WITH THE SOFTWARE OR THE USE OR OTHER DEALINGS IN
%	THE SOFTWARE.
%
%
%-----------------------------------------------------------------------------------------------------------------------------------------------%


%============================================================================%
%
%	DOCUMENT DEFINITION
%
%============================================================================%

%we use article class because we want to fully customize the page and dont use a cv template
\documentclass[10pt,A4]{article}


%----------------------------------------------------------------------------------------
%	ENCODING
%----------------------------------------------------------------------------------------

%we use utf8 since we want to build from any machine
\usepackage[utf8]{inputenc}

%----------------------------------------------------------------------------------------
%	LOGIC
%----------------------------------------------------------------------------------------

% provides \isempty test
\usepackage{xifthen}

%----------------------------------------------------------------------------------------
%	FONT
%----------------------------------------------------------------------------------------

% some tex-live fonts - choose your own

%\usepackage[defaultsans]{droidsans}
% \usepackage[default]{comfortaa}
% \usepackage{cmbright}
% \usepackage[default]{raleway}
%\usepackage{fetamont}
\usepackage{gillius}
% \usepackage[light,math]{iwona}
%\usepackage[sfdefault]{roboto}
\usepackage{lmodern}

% set font default
\renewcommand*\familydefault{\sfdefault}
\usepackage[T1]{fontenc}

% more font size definitions
\usepackage{moresize}

\usepackage{fontawesome7}

%----------------------------------------------------------------------------------------
%	PAGE LAYOUT  DEFINITIONS
%----------------------------------------------------------------------------------------

%debug page outer frames
%\usepackage{showframe}
\usepackage{paracol}

%define page styles using geometry
\usepackage[a4paper]{geometry}

% for example, change the margins to 2 inches all round
\geometry{top=1cm, bottom=0cm, left=0.5cm, right=0cm,}


%less space between header and content
\setlength{\headheight}{-5pt}


%customize entries left, center and right
%\lhead{}
%\chead{ \small{Jan Küster  $\cdot$ Consultant and Software Engineer $\cdot$  Bremen, Germany  $\cdot$  \textcolor{sectcol}{\textbf{info@jankuester.com}}  $\cdot$ +49 176 313 *** **}}
%\rhead{}


%indentation is zero
\setlength{\parindent}{0mm}

%----------------------------------------------------------------------------------------
%	TABLE /ARRAY DEFINITIONS
%----------------------------------------------------------------------------------------

%for layouting tables
\usepackage{multicol}
\usepackage{multirow}

%extended aligning of tabular cells
\usepackage{array}

\newcolumntype{x}[1]{%
>{\raggedleft\hspace{0pt}}p{#1}}%


%----------------------------------------------------------------------------------------
%	GRAPHICS DEFINITIONS
%----------------------------------------------------------------------------------------

%for header image
\usepackage{graphicx}

%for floating figures
\usepackage{wrapfig}
\usepackage{float}
%\floatstyle{boxed}
%\restylefloat{figure}

%for drawing graphics
\usepackage{tikz}
\usetikzlibrary{shapes, backgrounds,mindmap, trees, decorations.pathreplacing}


%----------------------------------------------------------------------------------------
%	Color DEFINITIONS
%----------------------------------------------------------------------------------------
\usepackage{transparent}
\usepackage{color}

%accent color
\definecolor{complcol}{RGB}{250,150,10}

%dark background color
\definecolor{bgcol}{RGB}{110,110,110}

%light background / accent color
\definecolor{softcol}{RGB}{225,225,225}

\definecolor{sectcol}{RGB}{0,120,150}


%Package for links, must be the last package used
\usepackage[hidelinks]{hyperref}

%============================================================================%
%
%
%	DEFINITIONS
%
%
%============================================================================%

%----------------------------------------------------------------------------------------
% 	ARROW GRAPHICS in Tikz
%----------------------------------------------------------------------------------------

% a six pointed arrow poiting to the left
\newcommand{\tzlarrow}{(0,0) -- (0.2,0) -- (0.3,0.2) -- (0.2,0.4) -- (0,0.4) -- (0.1,0.2) -- cycle;}

% include the left arrow into a tikz picture
% param1: fill color
%
\newcommand{\larrow}[1]
{\begin{tikzpicture}[scale=0.58]
	 \filldraw[fill=#1!100,draw=#1!100!black]  \tzlarrow
 \end{tikzpicture}
}

\newcommand{\tzlbrace}[1]
{
	\begin{tikzpicture}[ultra thick,scale=0.5,baseline={(current bounding box.base)},
	yshift=-7pt]
		\draw [decorate,#1,
			decoration = {brace}] (0,0) --  (0,1);
		\draw [decorate,#1,
			decoration = {brace}] (0.3,0) --  (0.3,1);
	\end{tikzpicture}
}
\newcommand{\tzrbrace}[1]
{
	\begin{tikzpicture}[ultra thick,scale=0.5,baseline={(current bounding box.base)},
	yshift=-7pt]
	\draw [decorate,#1,
		decoration = {brace,mirror}] (0,0) --  (0,1);
	\draw [decorate,#1,
		decoration = {brace,mirror}] (0.3,0) --  (0.3,1);
	\end{tikzpicture}
}

% a six pointed arrow poiting to the right
\newcommand{\tzrarrow}{ (0,0.2) -- (0.1,0) -- (0.3,0) -- (0.2,0.2) -- (0.3,0.4) -- (0.1,0.4) -- cycle;}

% include the right arrow into a tikz picture
% param1: fill color
%
\newcommand{\rarrow}
{\begin{tikzpicture}[scale=0.7]
	\filldraw[fill=sectcol!100,draw=sectcol!100!black] \tzrarrow
 \end{tikzpicture}
}

%----------------------------------------------------------------------------------------
%	custom sections
%----------------------------------------------------------------------------------------

% create a coloured box with arrow and title as cv section headline
% param 1: section title
%
\newcommand{\cvsection}[1]
{
\colorbox{sectcol}{\mystrut \makebox[\linewidth][l]{
\larrow{bgcol} \hspace{-8pt} \larrow{bgcol} \hspace{-8pt} \larrow{bgcol} \textbf{\textcolor{white}{\uppercase{#1}}}\hspace{4pt}
}}\\
}

% create a coloured arrow with title as cv meta section section
% param 1: meta section title
%
\newenvironment{metasection}[1] {
    \hspace{10pt}
	\vspace{6pt}
	\begin{center}
		\textcolor{white}{\large{\uppercase{#1}}}\\
	\normalsize
	\parbox{0.7\linewidth}{\textcolor{white}	\hrule}
    }{
    \end{center}
    \hspace{10pt}
}

%----------------------------------------------------------------------------------------
%	 CV EVENT
%----------------------------------------------------------------------------------------

% creates a stretched box as cv entry headline followed by some paragraphs about
% the work you did
% param 1:	event time i.e. 2014 or 2011-2014 etc.
% param 2:	event name (what did you do?)
% param 3:	institution (where did you work / study)
% param 4:	list of paragraphs outlining your contributions, see example usage below
%
\newcommand{\cvevent}[5]
{
    \hspace{3pt}
    \parbox{1\linewidth-3pt}{
    \vspace{6pt}
    \sbox0{\textcolor{black}{\textbf{#2}}}%
    \sbox1{\textcolor{bgcol}{\textcolor{complcol}{#3}\ifthenelse{\equal{#4}{}}{}{ (#4)}, #1}}%
    \ifdim\dimexpr\wd0+\wd1+1em\relax>\linewidth
        \noindent\usebox0\par\noindent\hspace*{\fill}\usebox1
    \else
        \noindent\usebox0\hfill\usebox1
    \fi
    \vspace{0pt}
    \textcolor{softcol}{\hrule}
    \vspace{6pt}
    \cvlist {#5}
    \vspace{-6pt}
    }
}

% formats a list of strings with variable length for use in `\cvevent`
% param 1: a list of strings outlining your contributions
\newcommand{\cvlist}[1] {
    \foreach \listitem in {#1}
    {
        \parbox{1\linewidth}{\larrow{softcol} \listitem\vspace{3pt}}
    }
}

% creates a stretched box as
\newcommand{\cveventmeta}[2]
{
	\mbox{\mystrut \hspace{87pt}\textit{#1}}\\
	#2
}

%----------------------------------------------------------------------------------------
% CUSTOM STRUT FOR EMPTY BOXES
%----------------------------------------- -----------------------------------------------
\newcommand{\mystrut}{\rule[-.3\baselineskip]{0pt}{\baselineskip}}


% use to vertically center content
% credits to: http://tex.stackexchange.com/questions/7219/how-to-vertically-center-two-images-next-to-each-other
\newcommand{\vcenteredinclude}[1]{\begingroup
\setbox0=\hbox{\includegraphics{#1}}%
\parbox{\wd0}{\box0}\endgroup}

% use to vertically center content
% credits to: http://tex.stackexchange.com/questions/7219/how-to-vertically-center-two-images-next-to-each-other
\newcommand*{\vcenteredhbox}[1]{\begingroup
\setbox0=\hbox{#1}\parbox{\wd0}{\box0}\endgroup}

%----------------------------------------------------------------------------------------
%	ICON-SET EMBEDDING
%----------------------------------------------------------------------------------------

% at this point we simplify our icon-embedding by simply referring to a set of png images.
% if you find a good way of including svg without conflicting with other packages you can
% replace this part
\newcommand{\icon}[3]{\makebox(#2, #2){\textcolor{#3}{\csname fa#1\endcsname}}}	%icon shortcut
\newcommand{\icontext}[4]{ 						%icon with text shortcut
	\vcenteredhbox{\icon{#1}{#2}{#4}} \vcenteredhbox{\textcolor{#4}{#3}}
}
\newcommand{\iconhref}[5]{ 						%icon with website url
    \vcenteredhbox{\icon{#1}{#2}{#5}} \href{#4}{\textcolor{#5}{#3}}
}
\newcommand{\iconemail}[5]{ 						%icon with email link
    \vcenteredhbox{\icon{#1}{#2}{#5}} \href{mailto:#4}{\textcolor{#5}{#3}}
}
\usepackage{fontspec}
\newfontfamily{\FA}[
	Path=fonts/,
]{fa-solid-900}
\newcommand{\iconnew}[3]{\makebox(#2, #2){\textcolor{#3}{\FA\symbol{"#1}}}}
\newcommand{\iconnewtext}[4]{ 						%icon with text shortcut
	\vcenteredhbox{\iconnew{#1}{#2}{#4}} \vcenteredhbox{\textcolor{#4}{#3}}
}
\newcommand{\iconnewhref}[5]{ 						%icon with website url
	\vcenteredhbox{\iconnew{#1}{#2}{#5}} \href{#4}{\textcolor{#5}{#3}}
}
\newcommand{\iconnewemail}[5]{ 						%icon with email link
	\vcenteredhbox{\iconnew{#1}{#2}{#5}} \href{mailto:Jérémie Libeau <#4>?subject=\emailSubject}{\textcolor{#5}{#3}}
}
\newcommand{\icontel}[5]{ 						%icon with email link
	\vcenteredhbox{\icon{#1}{#2}{#5}} \href{tel:#3}{\textcolor{#5}{#4}}
}
\newcommand{\iconcustomtext}[4]{
	\vcenteredhbox{\makebox(#2, #2){\includegraphics[width=#2px]{#1}}} \vcenteredhbox{\textcolor{#4}{#3}}
}

%============================================================================%
%
%
%
%	DOCUMENT CONTENT
%
%
%
%============================================================================%
